\subsection{纯量限制}

命题 7. 设 A、B 是两个环， \(\rho\) 是 A 到 B 的同态。设 E 是忠实平坦右 B-模。为使 \(\rho^{*}(\mathrm{E})\) 是平坦（相应地忠实平坦）的右 A-模，必须且只需 B 是平坦（相应地忠实平坦）的右 A-模。

把第 2 目的命题 4 应用于以 B、A、E、B 分别代替 R、S、E、F 的情形（B 上的右 A-模结构由 \(\rho\) 定义），可以看出

B 是平坦（相应地忠实平坦）A-模，当且仅当 \(E \otimes_{B} B = \rho^{*}(E)\) 是平坦（相应地忠实平坦）A-模。

\textbf{注}\par

\begin{enumerate}
\def\labelenumi{(\arabic{enumi})}
\item
  命题 7 表明，为使 B 是忠实平坦 A-模，只需存在一个忠实平坦 B-模，它同时是忠实平坦 A-模。
\item
  设 A、B、C 是三个环， \(\rho: A \rightarrow B\) 与 \(\sigma: B \rightarrow C\) 是两个环同态。命题 7 表明：若 C 是忠实平坦 B-模且 B 是忠实平坦 A-模，则 C 是忠实平坦 A-模。若 C 是忠实平坦 B-模又是忠实平坦 A-模，则 B 是忠实平坦 A-模（为确定起见，把模都取为右模）。另一方面，B 与 C 可以都是忠实平坦 A-模而 C 不是忠实平坦 B-模（习题 7）。
\end{enumerate}

